\section{Worked example: a multimodal grid}
\label{ch:worked}

The theory, data contract, and software workflow have so far been presented
separately. The grid reunites them in one reproducible calculation. Its small
network makes the route and modal mechanisms visible while exercising the same
validation, welfare, decomposition, and plotting commands used by larger
applications.

\subsection{Economic setting}

The worked example uses the economic-geography model. Transit is a second
mode in its transport network; the example does not use the urban
residence--workplace model. It has \GridNodeCount{} economic locations
arranged on a five-by-five grid. Adjacent locations are connected by
reciprocal road arcs. Transit is also available on the four links along the
central east--west axis and has a larger baseline modal share there. The
committed mode shifters rationalize these baseline shares; \(\eta\) governs
their response to a cost change rather than identifying relative cost levels
from shares alone. Each transit stop has an explicit terminal identifier, so
the example can include congestion from traffic using the two endpoint
terminals of a trip. Figure
\ref{fig:grid-network} shows the network.

\begin{figure}[H]
\centering
\includegraphics[width=0.60\textwidth]{grid-network.pdf}
\caption{Synthetic road grid and transit spine. Thin gray links are roads;
thick blue links also carry transit. Gold node area represents economic
activity, and blue rings identify transit terminals. Road and transit
connections are reciprocal.}
\label{fig:grid-network}
\end{figure}

Activity is positive throughout the grid and is higher near the central spine,
the north--south cross-axis, and the central node. A small eastward gradient
breaks otherwise arbitrary ranking ties. The input builder assigns symmetric
directional flows and generates location margins from the same accounting
system. The package therefore reconstructs the route kernel without silently
balancing the network.

To keep the equilibrium forces comparable with the paper while making the
network fully reproducible, the example uses
\[
 \alpha=0.10,\qquad \beta=-0.30,\qquad \sigma=9,
 \qquad \eta=1.099.
\]
These reference values mirror the paper calibration described in Section
\ref{ch:data}, drawing on
\citet{AllenArkolakis2014,AllenArkolakis2020} and
\citet{FuchsWong2026Multimodal}.
Modes enter through \code{ChoiceLogsum}. Road congestion has elasticity
0.08, and the illustrative endpoint-terminal elasticity for transit is 0.025.
The two congestion values make the model comparisons visible; they are not
estimates for a particular city.

\subsection{Build the inputs, run the model, and draw the results}

The committed input has \GridDirectedRoadArcCount{} directed road arcs,
\GridPhysicalRoadLinkCount{} reciprocal road links, and
\GridActiveEdgeModeCount{} positive-flow edge--mode pairs. The difference
between the first and third counts is the eight directed transit edge--mode
rows on the central spine.

The input builder is
\code{examples/grid\_multimodal/build\_inputs.jl}. It first assigns positive
labor and income across the grid. It then writes reciprocal horizontal
and vertical road arcs and adds transit rows only along the central axis. The
following excerpt shows the construction that an empirical preprocessing
script must replace:
\begin{lstlisting}[language=Julia]
const GRID_SIZE = 5
const TRANSIT_ROW = 3

node_id(row, column) = @sprintf("r%dc%d", row, column)

for row in 1:GRID_SIZE, column in 1:(GRID_SIZE - 1)
    west = node_id(row, column)
    east = node_id(row, column + 1)
    link = @sprintf("H_r%d_c%d", row, column)
    road = road_flow(:horizontal, row, column)
    edge_row(io, link * "_E", link, west, east, "road", road)
    edge_row(io, link * "_W", link, east, west, "road", road)
    if row == TRANSIT_ROW
        transit = transit_flow(column)
        edge_row(io, link * "_E", link, west, east,
                 "transit", transit;
                 origin_terminal=terminal_id(row, column),
                 destination_terminal=terminal_id(row, column + 1))
        edge_row(io, link * "_W", link, east, west,
                 "transit", transit;
                 origin_terminal=terminal_id(row, column + 1),
                 destination_terminal=terminal_id(row, column))
    end
end
\end{lstlisting}
The complete script also writes vertical links and explicit origin and
destination terminal identifiers. Symmetric directional flows make incoming
and outgoing traffic equal at each node. This is a transparent device for a
synthetic example, not a balancing rule that should be imposed on empirical
commuting or freight data.

Rebuild the CSVs and inspect the first rows before running the model:
\begin{lstlisting}[language=bash]
julia examples/grid_multimodal/build_inputs.jl
head examples/grid_multimodal/data/nodes.csv
head examples/grid_multimodal/data/edge_modes.csv
\end{lstlisting}
The configuration names the input files, selects
\code{economic\_geography} and \code{ChoiceLogsum}, declares road and terminal
congestion, and defines a one-percent bidirectional road policy. These are the
fields an analyst changes after replacing the CSVs:
\begin{lstlisting}
[model]
spatial_specification = "economic_geography"
alpha = 0.10
beta = -0.30
sigma = 9.0
eta = 1.099
modal_specification = "choice_logsum"

[policy]
mode = "road"
unit = "both"
shock_fraction = 0.01
\end{lstlisting}

Validation should precede the numerical work. From the repository root, run:
\begin{lstlisting}[language=bash]
julia --project=. bin/tnw.jl validate \
  examples/grid_multimodal/config.toml
julia --project=. bin/tnw.jl analyze \
  examples/grid_multimodal/config.toml
julia --project=. bin/tnw.jl decompose \
  examples/grid_multimodal/config.toml
julia --project=. bin/tnw.jl sensitivity \
  examples/grid_multimodal/config.toml
\end{lstlisting}
The analysis writes directed and physical-link welfare tables. The
decomposition adds the closure and component columns used below, while the
sensitivity command writes one row for each declared parameter value. Each run
also updates \code{diagnostics.json} and \code{run\_manifest.json}.

The figures, tables, and TeX macros in this section are generated from those
validated outputs rather than transcribed by hand. The documentation asset
target rebuilds the inputs, reruns the example, and redraws the figures:
\begin{lstlisting}[language=bash]
make practitioner-guide-assets
\end{lstlisting}
For an external application, the equivalent standalone map command is given in
Subsection \ref{sec:own-data-maps}; it reads
\code{welfare\_physical.csv} and can use line and polygon GeoJSON layers.

The policy improves one physical road link at a time by lowering primitive
cost by one percent in both directions. Transit remains available throughout
the equilibrium response. The result therefore measures how a road
improvement affects a network in which users can also change mode.

\subsection{Baseline checks}

The reconstructed route kernel has spectral radius
\(\GridRouteSpectralRadius\), safely below one. The full equilibrium Jacobian
has condition number \(\GridConditionF\), below the package gate of
\(10^{12}\). The maximum inverse-gap residual is
\(\GridInverseGapError\), the closure-ladder residual is
\(\GridLadderError\), and the largest analytical component-reconstruction
residual is \(\GridChannelError\). These checks establish that the reported
derivatives solve the stated numerical model. They do not turn the synthetic
parameters into empirical estimates.

\subsection{Traffic and primitive-cost rankings}

The mean Traditional-approach elasticity is \GridMeanTraditional, compared with a mean
primitive-cost welfare elasticity of \GridMeanExtended. The median and maximum
primitive-cost elasticities are \GridMedianExtended{} and
\GridMaximumExtended. For the configured \GridPolicyShockPercent{} percent
primitive improvement, the mean primitive-cost elasticity implies a
first-order log-welfare gain of
\GridMeanExtendedLogGain, approximately \GridMeanExtendedPercentGain{} percent.
This is a local benefit prediction, not a project net present value.

\begin{figure}[H]
\centering
\includegraphics[width=0.70\textwidth]{grid-welfare-map.pdf}
\caption{Extended-approach primitive-cost welfare elasticities for physical
road-link policies. Multiplying an elasticity by \(0.01\) gives the
first-order welfare gain from the configured one-percent improvement. The blue
line and terminal rings mark links on which transit is also available.}
\label{fig:grid-welfare}
\end{figure}

Benefits are largest on the central row and column, where economic activity
and network use are concentrated. The transit spine does not make the central
road links unimportant. It changes how their direct cost savings propagate
through modal allocation and congestion.

\begin{figure}[H]
\centering
\includegraphics[width=0.66\textwidth]{grid-scatter.pdf}
\caption{Traditional- and Extended-approach welfare elasticities for the
\GridPhysicalRoadLinkCount{} physical road improvements. The dashed line is
the 45-degree line. Labels identify the largest upward and downward rank
changes.}
\label{fig:grid-scatter}
\end{figure}

The Pearson correlation is \GridPearson, and the Spearman correlation is
\GridSpearman. The close linear relation reflects the importance of baseline
traffic. It does not imply identical rankings. Link \GridRiseLink{} rises from
\GridRiseTraditionalRank{} under the Traditional approach to
\GridRiseExtendedRank{} under the primitive-cost effect. Link \GridFallLink{}
falls from \GridFallTraditionalRank{} to \GridFallExtendedRank. These links
have similar traffic but differ in their access to the high-activity cross-axis
and the transit spine. The Extended approach values that difference in
equilibrium incidence.

\begin{table}[H]
\centering
\caption{Highest-ranked physical road links under the primitive-cost welfare
effect}
\label{tab:grid-top-links}
% Generated by plots/example_assets.py; do not edit.
\begin{tabular}{lrrrr}
\toprule
Link & Traditional & Extended & Traditional rank & Extended rank \\
\midrule
Row 3, columns 3-4 & 0.021040 & 0.010177 & 1 & 1 \\
Row 3, columns 2-3 & 0.020960 & 0.010069 & 2 & 2 \\
Row 3, columns 4-5 & 0.019120 & 0.009481 & 3 & 3 \\
Row 3, columns 1-2 & 0.018880 & 0.009214 & 4 & 4 \\
Column 3, rows 3-4 & 0.015640 & 0.007433 & 5 & 5 \\
Column 3, rows 2-3 & 0.015240 & 0.007312 & 6 & 6 \\
Column 3, rows 4-5 & 0.014440 & 0.007294 & 7 & 7 \\
Column 3, rows 1-2 & 0.013240 & 0.006849 & 8 & 8 \\
\bottomrule
\end{tabular}

\end{table}

The highest-ranked link, \GridTopLink, has a Traditional-approach elasticity of
\GridTopTraditional, a full realized-cost effect of \GridTopRealized, and a
primitive-cost effect of \GridTopPrimitive. The highest-ranked links are
stable across the traffic and primitive-cost measures. Ranking changes
occur farther down the list, where direct traffic differences are smaller and
equilibrium propagation has more scope to matter. A high correlation or
agreement at the top therefore does not establish that the two criteria are
equivalent.

\subsection{Closure ladder}

The common-baseline ladder separates the modeled effects without recalibrating the
starting equilibrium. Applying the cases defined in Section
\ref{ch:decomposition}, the mean realized-cost elasticities are
\[
\begin{array}{c|ccccc}
  & NC & NT & F & FM & FR\\
\hline
\text{mean elasticity}
  & \GridMeanRealizedNC
  & \GridMeanRealizedNT
  & \GridMeanRealizedF
  & \GridMeanRealizedFM
  & \GridMeanRealizedFR .
\end{array}
\]
Moving from \(NC\) to \(NT\) adds road congestion; moving from \(NT\) to \(F\)
adds terminal congestion. Comparing \(F\) with \(FM\) or \(FR\) isolates the
local mode and route wedges. Their mean normalized values are
\GridMeanModeWedge{} and \GridMeanRouteWedge, respectively. Because these are
normalized multipliers, their signs should be interpreted together with
\(\rho\), not as welfare levels on their own.

\begin{figure}[H]
\centering
\includegraphics[width=0.98\textwidth]{grid-decomposition.pdf}
\caption{Closure ladder and exact welfare decomposition. Panel (a) moves from
the Traditional approach through the no-congestion,
road-congestion, and full realized-cost cases to the Extended-approach
primitive-cost effect. The hollow
\(FM\) and \(FR\) markers hold modes and routes fixed. Panel (b) reports the
signed contributions to Traditional-approach welfare minus primitive-cost welfare;
positive values widen the gap and negative values offset it.}
\label{fig:grid-decomposition}
\end{figure}

Panel (a) of Figure \ref{fig:grid-decomposition} shows which additions change
the mean welfare elasticity. Panel (b) accounts for the gap between the
Traditional approach and the primitive-cost mean. The common externality term is
\GridExternalityComponent, but general-equilibrium propagation contributes
\GridPropagationComponent{} and nearly cancels it. Road and terminal
congestion contribute \GridRoadComponent{} and \GridTerminalComponent.
Primitive-cost pass-through contributes \GridPassThroughComponent{} in the
opposite direction. The components sum exactly to the mean difference between
the Traditional and Extended approaches.

The cancellation reflects opposing economic effects rather than a failed
identity: the reported residuals are below the stated tolerance. Spatial
externalities generate a common scale, whereas market-access propagation
depends on each link's position in the network, so the two contributions can
have opposite signs. The road and terminal terms measure changes in congestion
exposure. Pass-through measures a separate margin: induced traffic partly
offsets the primitive cost reduction before it reaches aggregate edge costs.
Reporting signed levels displays these offsets without dividing by a net gap
that is small relative to its components.

\subsection{Adapting the example}

First reproduce the committed grid before replacing its inputs. Section
\ref{sec:own-data-workflow} then gives the complete workflow for creating an
external data project, updating the package without overwriting that project,
validating the new baseline, and recording a new result vintage. The generic
loader requires model-ready flows and does not automatically balance,
symmetrize, or rescale them.

Other examples exercise specialized parsers and geometries and are catalogued
in the repository's examples index. Subsection \ref{ch:diagnostics} defines
the distinction between numerical verification and external validation that
also governs those examples.
